\documentclass[10pt,a4paper]{amsart}
\usepackage[margin=19mm]{geometry}
\usepackage{amsmath,amssymb,mathtools}
\usepackage[scr=rsfs,cal=euler]{mathalfa}
\usepackage{xfrac}
\usepackage[unicode,hidelinks,bookmarks=false]{hyperref}
\newcommand{\ie}{i.e.\ }
\newcommand{\cf}{cf.\ }
\newcommand{\resp}{respectively\ }
\newcommand{\Subsection}{\subsection}
\providecommand{\nobreakdash}{\nobreak\hbox{-}}
\setlength{\emergencystretch}{2em}
\multlinegap0pt
% ---------------------------------------------------------------------------
% Editorial AMS wrapper prelude. The theorem environments and the mathematical macros below are
% transcribed from the paper's own preamble (cached-originals/source0.tex lines 9--60); the AMS amsart
% class, the named format packages of the pinned TeX Live runtime and this editorial formatting are not
% part of the author's text. No mathematics is changed.
% ---------------------------------------------------------------------------
\DeclareMathOperator{\spn}{span}
\DeclareMathOperator{\ind}{ind}
\DeclareMathOperator{\tr}{tr}
\DeclareMathOperator{\dom}{dom}
\def\arctg{{\rm arctg\,}}
\def\d{\displaystyle}
\def\t#1{$^{\bf *}$\bf #1{.}}
\newtheorem{thm}{Theorem}[section]
\newtheorem{prop}[thm]{Proposition}
\newtheorem{lem}[thm]{Lemma}
 \newtheorem{rem}[thm]{Remark}
 \newtheorem{cor}[thm]{Corollary}
\newtheorem{defn}[thm]{Definition}
\newtheorem{ex}[thm]{Example}
\DeclareMathOperator{\supp}{supp}
\DeclareMathOperator{\ext}{ext}
\newcommand{\C}{\mathbb{C}}
\newcommand{\R}{\mathbb{R}}
\renewcommand{\S}{\mathbb{S}}
\newcommand{\N}{\mathbb{N}}
\newcommand{\T}{\mathbb{T}}
\newcommand{\M}{\mathbb{M}}
\newcommand{\D}{\mathbb{D}}
\newcommand{\Z}{\mathbb{Z}}
\newcommand{\la}{\lambda}
\def\labelenumi{\bf \arabic{enumi}.}
\def\tlab{$\bf *$\labelenumi}
\long\def\comment#1{{}}
\def\ph{\phantom{-}}
\def\p{\ph}

\begin{document}
\title{Indeterminate Jacobi operators II}
\author{Christian Berg and Ryszard Szwarc}
\date{Source: arXiv:2510.04690v2; distributed under CC BY 4.0}
\maketitle
\noindent\emph{Source, version and licence.} \emph{Indeterminate Jacobi operators II}, by Christian Berg and Ryszard Szwarc. The source of this editable unit is the e-print of \textbf{version v2} of arXiv:2510.04690, https://arxiv.org/abs/2510.04690v2, whose material is distributed under the Creative Commons Attribution 4.0 International licence, http://creativecommons.org/licenses/by/4.0/, as recorded in the arXiv licence metadata for this paper and shown on that abstract page. The whole of the body below is version v2, the newest version of the paper. Full rights notice: licenses/arxiv-2510.04690-CC-BY-4.0.txt.\par\smallskip
\noindent\textit{Editorial scope.} Counted unit: original Proposition 3.1 of the article, the statement carried by the source environment \texttt{prop} with source label \texttt{thm:1-2}, cached-originals/source0.tex lines 355--363, printed as \textbf{Proposition 3.1} exactly as in the article, and described there as one of the two results the authors believe to be new. It is delivered together with the author\'s own complete proof as the single primary proof body, source0.tex lines 365--381 (17 lines): the \texttt{proof} environment that immediately follows the statement and proves parts (i), (ii) and (iii) in turn. No proof marker is added and no mathematics is written, repaired or re-derived; the author\'s notation and the printed slips of the source are kept literal. Necessary complete context is retained uncounted, in source order: source0.tex lines 64--71, the defining equations of the Jacobi operator and its indeterminate setting; lines 122--124, the identity $B(u)+tD(u)$ used in the counted proof; lines 165--176, Theorem 1.2 on the solutions of the second kind; lines 198--200 and 215--217, the support identity and the description of the N-extremal measures; lines 222--223, Theorem 1.3, the main result of the paper that the counted proposition serves; lines 261--278, Proposition 2.1 with the vanishing integrals of the polynomial approximants; lines 304--312, Proposition 2.2 parametrizing the N-extremal measures; lines 322--323, Proposition 2.3 on the countable support of those measures; and lines 383--389, Corollary 3.2, the companion limit statement. The article\'s own numbering is reproduced by explicit counter settings: the counter \texttt{thm}, declared by \texttt{\textbackslash newtheorem\{thm\}\{Theorem\}[section]} and shared by prop, lem, rem, cor, defn and ex, is set before each retained passage, so the counted statement prints as Proposition 3.1, the main theorem as Theorem 1.3, the companion as Corollary 3.2 and the quoted results as Propositions 2.1, 2.2 and 2.3, exactly as in the article. Every \texttt{\textbackslash ref} and \texttt{\textbackslash eqref} inside every retained passage resolves inside this delivered file, and the single citation that occurs in the retained passages, \texttt{Sch}, is delivered as the author\'s own bibliography entry transcribed verbatim from the article\'s own \texttt{thebibliography} with the author\'s own printed number. The retained passages contain no figure environment and no graphics-inclusion command, so no artwork is redistributed. The source\'s own class is the standard \texttt{article}; the e-print ships no class or style file, so no publisher class is shipped, used or redistributed, and the wrapper is the AMS \texttt{amsart} class with the paper\'s own macros and theorem declarations transcribed from its preamble. Every declared body of this unit is an exact, unmodified span of source0.tex: no body is a bounded passage comparison, \texttt{omit} was not used anywhere, and consequently no fidelity receipt is asserted in delivery-evidence.json. The complete concatenated original source is bundled as sources/186462/source0.tex. Omitted from this editable unit and therefore absent from it are source0.tex lines 1--63; 72--121; 125--164; 177--197; 201--214; 218--221; 224--260; 279--303; 313--321; 324--354; 364--364; 382--382; 390--512; 516--535, that is the article\'s own preamble, title block and abstract, the citation-dense introduction, the proofs of Theorems 1.2 and 1.3, the remaining parts of the paper, and the remaining bibliography entries, all of which remain present in the bundled source but are not delivered here. The AMS \texttt{amsart} wrapper, the named format packages of the pinned TeX Live runtime and this editorial source, licence and scope paragraph are editorial formatting only.\par\smallskip
\setcounter{section}{1}
\setcounter{equation}{1}
\setcounter{thm}{0}
\begin{equation}\label{eq:Jac}
J=\begin{pmatrix}
b_0 & a_0 & 0 & \hdots\\
a_0 & b_1 & a_1 & \hdots\\
0 & a_1 & b_2 & \hdots\\
\vdots &\vdots & \vdots & \ddots
\end{pmatrix},
\end{equation}

\setcounter{section}{1}
\setcounter{equation}{4}
\setcounter{thm}{1}
\begin{equation}\label{eq:FGc}
F_c(z)=\sum_{n=0}^\infty c_np_n(z),\quad G_c(z)=\sum_{n=0}^\infty c_nq_n(z),\quad z\in\C.
\end{equation}

\setcounter{section}{1}
\setcounter{equation}{14}
\setcounter{thm}{1}
\begin{thm}\label{thm:pqmain} For all $z\in\C$ we have $\mathfrak{p}_z, \mathfrak{q}_z\notin D(T)$.

Let $u,v\in\C$ be given.
\begin{enumerate}
\item[(i)] There exists $\alpha\in\C$ such that $\mathfrak{p}_u+\alpha \mathfrak{p}_v\in D(T)$ if and only if $D(u,v)=0$. In the affirmative case $\alpha$  is uniquely determined as
$\alpha=B(u,v)$. 

\item[(ii)] There exists $\beta\in\C$ such that $\mathfrak{q}_u+\beta \mathfrak{q}_v\in D(T)$ if and only if $A(u,v)=0$. In the affirmative case $\beta$  is uniquely determined as $\beta=-C(u,v)$.

\item[(iii)] There exists $\gamma\in\C$ such that $\mathfrak{p}_u+\gamma \mathfrak{q}_v\in D(T)$ if and only if $B(u,v)=0$. In the affirmative case $\gamma$  is uniquely determined as $\gamma=-D(u,v)$. In particular $\mathfrak{p}_u+\gamma\mathfrak{q}_u\notin D(T)$ for all $u, \gamma\in\C$.
\end{enumerate}
\end{thm}

\setcounter{section}{1}
\setcounter{equation}{17}
\setcounter{thm}{2}
\begin{equation}\label{eq:supp}
\supp(\mu[v_0])=\{u\in\R \mid D(u,v_0)=0\},
\end{equation}

\setcounter{section}{1}
\setcounter{equation}{20}
\setcounter{thm}{2}
\begin{equation}\label{eq:N-ext2}
\{u\in\R \mid B(u,v_0)=0\}=\supp(\mu_s),
\end{equation}

\setcounter{section}{1}
\setcounter{equation}{21}
\setcounter{thm}{2}
  \begin{thm}\label{thm:concrete} For each $v_0\in\R$ the subspaces $P(v_0)$ and $M(v_0)$ are dense in $\ell^2$, while $Q(v_0)$ is dense in $e_0^\perp$.
\end{thm}

\setcounter{section}{2}
\setcounter{equation}{0}
\setcounter{thm}{0}
\begin{prop}\cite[Proposition 5.24]{Sch}\label{thm:An-Dn} For $u,v\in\C$ and $n\ge 0$ we have
\begin{eqnarray*}
A_n(u,v)&:=&(u-v)\sum_{k=0}^nq_k(u)q_k(v)=
a_n\left|\begin{array}{cc}
 q_{n+1}(u)&\;q_{n+1}(v)\\q_{n}(u)&\;q_{n}(v)\end{array}\right|,\\
B_n(u,v)&:=&-1+(u-v)\sum_{k=0}^n p_k(u)q_k(v)=
a_n\left|\begin{array}{cc}
p_{n+1}(u)&\;q_{n+1}(v)\\p_{n}(u)&\;q_{n}(v)\end{array}\right|,\\
C_n(u,v)&:=&1+(u-v)\sum_{k=0}^n q_k(u)p_k(v)=
a_n\left|\begin{array}{cc}
q_{n+1}(u)&\;p_{n+1}(v)\\
q_{n}(u)&\;p_{n}(v)\end{array}\right|,\\
D_n(u,v)&:=&(u-v)\sum_{k=0}^np_k(u)p_k(v)=
a_n\left|\begin{array}{cc}
p_{n+1(}u)&\;p_{n+1}(v)\\
p_{n}(u)&\;p_{n}(v)\end{array}\right|.
\end{eqnarray*}
\end{prop}

\setcounter{section}{2}
\setcounter{equation}{3}
\setcounter{thm}{1}
\begin{prop}\label{thm:supportNext}
\begin{enumerate}
\item[(i)] The  solution $\mu_t$ is a discrete measure with support equal to the countable zero set $\Lambda_t$ of the entire function $B(z)+tD(z)$, with the convention that $\Lambda_\infty$ is the zero set of $D$. We have $\Lambda_t\subset\R$ for $t\in\R^*$ and
the mass of $\mu_t$ at $u\in\Lambda_t$ equals $||\mathfrak{p}_u||^{-2}$. 
\item[(ii)] The supports of two different N-extremal solutions are disjoint and interlacing. Each point $x_0\in\R$ belongs to the support of a unique N-extremal measure $\mu_t$, where
$t\in\R^*$ is given  as $t=-B(x_0)/D(x_0)$ if $D(x_0)\neq 0$ and $t=\infty$ if $D(x_0)=0$.
\item[(iii)] If $x_0\in\R$ belongs to the support of the N-extremal measure $\mu_t$, then the measure $\mu_t-\mu_t(x_0)\delta_{x_0}$ is determinate.
\end{enumerate}
\end{prop}

\setcounter{section}{2}
\setcounter{equation}{4}
\setcounter{thm}{2}
\begin{prop}\label{thm:ONB} Let $(u_n)_{n\ge 1}$ be the points of the countable set $\Lambda_t$ supporting $\mu_t$, where $t\in\R^*$. The vectors $\mathfrak{p}_{u_n}/||\mathfrak{p}_{u_n}||, n\ge 1$ form an orthonormal basis for $\ell^2$.
\end{prop}

\setcounter{section}{3}
\setcounter{equation}{0}
\setcounter{thm}{0}
\begin{prop}\label{thm:1-2} (i) For $\mu\in V, v\in\C$ and any polynomial $p$ the functions $p(u)B(u,v), p(u)D(u,v)$ belong to $L^1(\mu)$ as functions of $u$ and
$$
\int p(u)B(u,v)\,d\mu(u)=\int p(u)D(u,v)\,d\mu(u)=0.
$$ 

(ii) For $v_0\in\R$ we have $B(u,v_0)\notin L^2(\mu)$ for any N-extremal measure $\mu$ different from the N-extremal measure $\mu_s$ supported by $\{u\in\R \mid B(u,v_0)=0\}$, cf. \eqref{eq:N-ext2}, while $B(u,v_0)$ is the zero element in $L^2(\mu_s)$.

(iii) For $v_0\in\R$ we have $D(u,v_0)\notin L^2(\mu)$ for any N-extremal measure $\mu\neq \mu[v_0]$, while $D(u,v_0)$ is the zero element in $L^2(\mu[v_0])$.
 \end{prop} 

\setcounter{section}{3}
\setcounter{equation}{0}
\setcounter{thm}{1}
\begin{proof} (i). With the notation of \eqref{eq:FGc} we have
\begin{eqnarray*}
p(u)B(u,v)&=&-p(u)+p(u)(u-v)F_{(q_k(v))}(u),\\
 p(u)D(u,v)&=&p(u)(u-v)F_{(p_k(v))}(u),
\end{eqnarray*}
so it is clear from the Cauchy-Schwarz inequality that $p(u)B(u,v), p(u)D(u,v)$ belong to $L^1(\mu)$ as functions of $u$ and furthermore that
$p(u)B_n(u,v)\to p(u)B(u,v)$ and $p(u)D_n(u,v)\to p(u)D(u,v)$  in $L^1(\mu)$ for $n\to\infty$. However,
\begin{equation}\label{eq:int0}
\int p(u)B_n(u,v)\,d\mu(u)=\int p(u)D_n(u,v)\,d\mu(u)=0
\end{equation}
 for $n>\deg(p)$ by Proposition~\ref{thm:An-Dn}, and (i) follows.

(ii). If $B(u,v_0)\in L^2(\mu)$ for an N-extremal measure $\mu$ it follows from \eqref{eq:int0} that $\int p(u)B(u,v_0)\,d\mu(u)=0$ for all polynomials $p$. Since the polynomials are dense in $L^2(\mu)$ we conclude that $B(u,v_0)=0$ for all $u\in\supp(\mu)$. Therefore $\mu$ must be the N-extremal measure supported by $\{u\in\R \mid B(u,v_0)=0\}$.

(iii). If $D(u,v_0)\in L^2(\mu)$ for an N-extremal measure $\mu$, we get that $D(u,v_0)=0$ for all $u\in\supp(\mu)$. This is not possible if $\mu\neq \mu[v_0]$ because then  
$\supp(\mu)$ is disjoint from $\supp(\mu[v_0])$ and this contradicts \eqref{eq:supp}, which also shows that  $D(u,v_0)$ is the zero element in $L^2(\mu[v_0])$.
\end{proof}

\setcounter{section}{3}
\setcounter{equation}{1}
\setcounter{thm}{1}
\begin{cor}\label{thm:limit} For any $v_0\in\R$ we have
\begin{eqnarray*}
&\lim_{n\to\infty} a_n^2[p_n(v_0)^2+p_{n+1}(v_0)^2]=\infty,\\
&\lim_{n\to\infty} a_n^2[q_n(v_0)^2+q_{n+1}(v_0)^2]=\infty,
\end{eqnarray*}
where $a_n$ is the entry $J_{n+1,n}$  of the Jacobi matrix \eqref{eq:Jac}.
\end{cor}

\begin{thebibliography}{99}
\bibitem[11]{Sch} K.~Schm\"{u}dgen, {\it The Moment Problem}. Graduate Texts in Mathematics Vol. 277. Springer International Publishing AG 2017. 
\end{thebibliography}
\end{document}
